\section{Relative comparison of the retained densities}
\label{cmp:section}

Closeness of the coefficients in an effective action does not by itself
compare the corresponding partition functions.  The difficulty here is
quite specific.  A retained density contains a signed covering sum, and
this sum can be small on configurations with many rough bonds.  Its
absolute difference is then an inadequate estimate for its relative
difference.

We prove the relative estimate by combining two facts.  Rough bonds are
rare under each of the two positive measures.  On the remaining
configurations, a local modification removes rough clusters at a
controlled cost.  An exact inclusion--exclusion identity then turns the
signed comparison into a convergent sum of trees.  Positivity will be
used for complete densities and complete covering sums, never for a
restricted signed sum.

\subsection{The finite-volume comparison problem}

Let
\[
 \Lambda_M=(\mathbb Z/M\mathbb Z)^2,\qquad v=|\Lambda_M|=M^2,
 \qquad q\in(S^3)^{\Lambda_M},
\]
where \(M\) is dyadic.  Write \(\dd q\) for product normalized Haar
measure, \(r_e=|q_x-q_y|\) for \(e=\langle xy\rangle\), and
\[
 D(q)=\{e:r_e>t_H\}.
\]
The large parameter \(H\) is the inverse coupling at the retained scale.
It is fixed throughout a comparison.  In the application,
\[
 t_H=H^{-1/2}p_H,\qquad p_H=(\log H)^{P_0},
\]
with the fixed exponent \(P_0\) chosen in the renormalization construction.

\begin{definition}[A mandatory covering sum]
\label{cmp:cover-definition}
A label \(\lambda\) consists of a finite site support \(P_\lambda\), a
load \(s_\lambda\ge1\), and a fixed nonempty inventory \(J_\lambda\) of
bonds whose endpoints belong to \(P_\lambda\).  Its weight
\(k_\lambda(q)\) is a bounded measurable function of \(q|_{P_\lambda}\);
every eligibility indicator is included in this function.  In
particular \(k_\lambda(q)=0\) unless \(J_\lambda\subset D(q)\).
The support contains every spin read by the function or by an
eligibility test, including a test asserting the absence of a flag.

A family is compatible when its supports are pairwise disjoint.  The
mandatory covering sum is
\begin{equation}
 \Xi(q)=
 \sum_{\substack{\Gamma\ {\rm compatible}\\
          \bigsqcup_{\lambda\in\Gamma}J_\lambda=D(q)}}
            \prod_{\lambda\in\Gamma}k_\lambda(q).
 \label{cmp:cover}
\end{equation}
The empty family has weight one.  Since inventories are nonempty,
\(\Xi(q)=1\) when \(D(q)=\varnothing\).
\end{definition}

The word mandatory refers to the equality of inventories in
\eqref{cmp:cover}.  In particular, \(\Xi\) is not the partition function
of an optional dilute gas.

Consider two pointwise positive densities
\begin{equation}
 \rho_i(q)=e^{X_i(q)}\Xi_i(q)>0,\qquad
 Z_i=\int\rho_i(q)\,\dd q,\qquad
 \dd\mu_i=Z_i^{-1}\rho_i\,\dd q,\quad i=1,2 .
 \label{cmp:densities}
\end{equation}
Use a common countable label set, putting unused weights equal to zero.
The functions in \eqref{cmp:densities} are the pointwise versions
provided by the exact blocking identity.

Here are the precise data needed for comparison.  A quantity described
as polynomial in \(H\) is bounded by \(CH^a\), with \(C,a\) fixed
independently of the period and the cutoff depths.

\begin{enumerate}
\item[\textup{(C1)}]
Each \(P_\lambda\) is contained in a connected lattice graph with at
most \(R_Hs_\lambda\) vertices, where \(R_H\) is polynomial in \(H\).
For a fixed \(A\ge16\),
\begin{align}
 \sup_x\sum_{\lambda:x\in P_\lambda}
          e^{As_\lambda}\|k_{i,\lambda}\|_\infty
       &\le w_H, \label{cmp:activity}\\
 \sup_x\sum_{\lambda:x\in P_\lambda}
          e^{As_\lambda}
          \|k_{2,\lambda}-k_{1,\lambda}\|_\infty
       &\le D_H\delta_K. \label{cmp:activity-difference}
\end{align}
Here \(w_H\) decreases faster than every inverse power of \(H\);
\(D_H<\infty\) is independent of the period and the cutoff depths.

\item[\textup{(C2)}]
If \(q,\widetilde q\) differ at \(n\) sites, then
\begin{align}
 |\log\rho_i(q)-\log\rho_i(\widetilde q)|
       &\le B_{\rho,H}n, \label{cmp:density-finite-energy}\\
 |X_i(q)-X_i(\widetilde q)|
       &\le B_{X,H}n, \label{cmp:exponent-finite-energy}
\end{align}
where \(B_{\rho,H},B_{X,H}\) are polynomial in \(H\).  Moreover
\begin{equation}
 \|X_2-X_1\|_\infty\le D_Hv\delta_K.
 \label{cmp:exponent-difference}
\end{equation}

\item[\textup{(C3)}]
For some \(c_*>0\) independent of \(H\),
\begin{equation}
 X_i(q)\le \epsilon_Hv-
                c_*H\sum_e\min(r_e^2,t_H^2),
 \qquad \epsilon_H=o(Ht_H^2).
 \label{cmp:stability}
\end{equation}
On the cap
\[
 {\cal C}_H=\{q:|q_x-\mathbf1|\le cH^{-1/2}
                         \text{ for every }x\},
\]
one has \(D(q)=\varnothing\) and
\begin{equation}
 X_i(q)\ge-\epsilon'_Hv,\qquad
 \epsilon'_H=o(Ht_H^2).
 \label{cmp:cap}
\end{equation}

\item[\textup{(C4)}]
Both positive measures are translation invariant and reflection
positive at the axis seams used to reflect lattice tiles.  The
reflection prescriptions are compatible with the tile events below.
Finally,
\[
 t_H\longrightarrow0,\qquad
 t_H^{-1}\le\operatorname{poly}(H),\qquad
 \frac{Ht_H^2}{(\log H)^2}\longrightarrow\infty .
 \label{cmp:thresholds}
\]
\end{enumerate}

All comparisons below are uniform over families satisfying these data.
The last subsection verifies the data for the retained densities of
Theorem~\ref{thm:rg-closure}.

\begin{theorem}[Relative comparison]
\label{thm:relative-comparison}
Suppose \textup{(C1)--(C4)} hold and
\(\delta_K\le C_H e^{-\gamma_H K}\), where \(\gamma_H>0\).
Fix \(0<\zeta_H<\gamma_H/2\).  For every sufficiently large fixed \(H\)
there are constants \(M_H,C_H,a_H>0\) such that, whenever \(M\ge M_H\),
there is a common measurable set \(G=G_{H,K,M}\) satisfying
\begin{align}
 \mu_i(G^c)&\le C_HM^2e^{-a_HK},\qquad i=1,2,
                    \label{cmp:bad-probability}\\
 \left|\frac{\Xi_2(q)}{\Xi_1(q)}-1\right|
   &\le \exp\!\left(C_HM^2\delta_K e^{\zeta_HK}\right)-1,
                       \qquad q\in G. \label{cmp:relative-cover}
\end{align}
The same estimate holds after interchanging \(1\) and \(2\).  Thus
\begin{equation}
 \left|\log\frac{\rho_2(q)}{\rho_1(q)}\right|
       \le C_HM^2\delta_K e^{\zeta_HK},\qquad q\in G.
 \label{cmp:relative-density}
\end{equation}
There are also \(c_H,d_H>0\) and \(K_H<\infty\) such that
\begin{equation}
 \left|\log\frac{Z_2}{Z_1}\right|\le C_He^{-d_HK}
 \quad\text{if}\quad
 K\ge K_H,\qquad M_H\le M,\qquad M^2\le e^{c_HK}.
 \label{cmp:partition-conclusion}
\end{equation}
\end{theorem}

\subsection{A joint rough-bond estimate under each positive law}

Choose a fixed constant \(C_1\), to be enlarged in the geometric
construction, and define the wider set of flagged bonds
\begin{equation}
 \widehat D(q)=\{e:r_e>t_H/C_1\}.
 \label{cmp:wider-flags}
\end{equation}
The distinction between \(D\) and \(\widehat D\) leaves room to modify a
configuration near its true rough bonds.

\begin{lemma}[Joint rarity]
\label{cmp:rarity}
For sufficiently large \(H\), there is
\(\sigma_H=cHt_H^2\), with \(c>0\) fixed after \(C_1\), such that
\begin{equation}
 \mu_i\{J\subset\widehat D\}\le e^{-\sigma_H|J|}
 \label{cmp:joint-rarity}
\end{equation}
for every set \(J\) of distinct bonds and \(i=1,2\).
\end{lemma}

\begin{proof}
Dropping the coverage and compatibility constraints only in an absolute
upper bound gives
\begin{equation}
 |\Xi_i(q)|
 \le\prod_\lambda(1+\|k_{i,\lambda}\|_\infty)
 \le\exp\!\left(\sum_\lambda\|k_{i,\lambda}\|_\infty\right)
 \le e^{vw_H}.
 \label{cmp:absolute-cover}
\end{equation}
All these sums converge by \eqref{cmp:activity}.  On \({\cal C}_H\) the
covering sum is exactly one.  A cap of radius \(cH^{-1/2}\) on \(S^3\)
has Haar measure at least \(c'H^{-3/2}\).  Consequently
\begin{equation}
 Z_i\ge
       \exp[-v(\epsilon'_H+C+\tfrac32\log H)].
 \label{cmp:denominator-cap}
\end{equation}

Place one specified flagged-bond event inside a \(2\)-by-\(2\) tile,
and reflect it into every tile.  The fully disseminated event requires
at least \(cv\) distinct flagged bonds.  On this event,
\(\sum_e\min(r_e^2,t_H^2)\ge c'vt_H^2\).  Equations
\eqref{cmp:stability}, \eqref{cmp:absolute-cover}, and
\eqref{cmp:denominator-cap} therefore imply
\begin{equation}
 \mu_i(\text{fully disseminated event})
 \le
 \exp\{-v[c''Ht_H^2-\epsilon_H-\epsilon'_H-w_H-C\log H]\}
 \le e^{-c'''vHt_H^2}.
 \label{cmp:disseminated}
\end{equation}

This is the classical chessboard estimate
\cite[Theorems 4.1--4.3]{FILS78}. To display the reflection hypotheses
used here, we include the following form of
successive Cauchy--Schwarz.  If \(f_z\) are nonnegative tile tests and
\(\theta_z f_z\) denotes their reflected placement in tile \(z\), then
\[
 \mu_i\!\left(\prod_z\theta_zf_z\right)
 \le
 \prod_z
 \mu_i\!\left(\prod_y\theta_y f_z\right)^{1/N_{\rm tile}}.
\]
At one seam, reflection positivity applies Cauchy--Schwarz to the two
half-torus functions.  It replaces them by their two reflected
duplicates and gives exponent \(1/2\) to each.  Repeat in each direction,
bisecting the dyadic array of tiles.  After all bisections a chosen tile
test has been reflected into every tile and has exponent
\(1/N_{\rm tile}\), which is the displayed inequality.

Partition \(J\) into finitely many orientation, parity, and tile-origin
classes.  One class contains at least a fixed fraction of its bonds,
with at most one specified test per tile.  Use the displayed inequality
with this test on those tiles and \(1\) elsewhere.  Applying
\eqref{cmp:disseminated} proves \eqref{cmp:joint-rarity}, after reducing
the constant in \(\sigma_H\).  This proof is performed separately for
each \(\mu_i\); it does not compare their normalizing constants.
\end{proof}

\subsection{Removing a rough cluster}

We first record the target-space fact used in the modification.

\begin{lemma}[Quantitative filling on \(S^3\)]
\label{cmp:filling}
An \(L_0\)-Lipschitz map from the boundary of the unit square to \(S^3\)
has an extension to the square with Lipschitz constant at most
\(C(1+L_0)^{3/2}\).  The extension can be chosen measurably as a function
of the boundary map.
\end{lemma}

\begin{proof}
Parameterize the boundary by a circle using a fixed bi-Lipschitz map
and write its image as \(g(\theta)\).  Its antipodal curve can be
covered by \(C(1+L_0/\eta)\) balls of radius \(2\eta\).  The total Haar
volume of these balls is at most
\(C(\eta^3+L_0\eta^2)\).  With
\(\eta=c(1+L_0)^{-1/2}\), this is smaller than the volume of \(S^3\).
Choose \(p\in S^3\) at distance at least \(c\eta\) from the antipodal
curve.  Then
\[
 F(r,\theta)=
 \frac{(1-r)p+rg(\theta)}{|(1-r)p+rg(\theta)|},
 \qquad 0\le r\le1,
\]
is well defined.  Indeed its denominator is bounded below by
\(c\eta\).  The radial derivative is \(O(\eta^{-1})\), and its angular
derivative is \(O(rL_0\eta^{-1})\).  In disk coordinates these imply the
claimed bound.  A fixed bi-Lipschitz map between disk and square changes
only the constant.

To make the choice measurable, fix a sufficiently fine finite net of
possible \(p\)'s and take the first one separated from the antipodal
curve by the smaller prescribed distance.  The volume argument with
slightly larger avoidance balls guarantees such a net point.
\end{proof}

\begin{lemma}[Disjoint editable regions]
\label{cmp:editing}
For fixed sufficiently large constants \(C_2,C_1\), in that order,
and \(M\ge M_H\), each configuration has a finite family of disjoint
editable site sets \(E_C\), with nonempty fixed inventories
\(D_C\subset D(q)\), such that
\[
 D(q)=\bigsqcup_C D_C.
\]
If \(n_C\) is the number of original flagged bonds assigned to this
region, then \(n_C\ge1\), the assigned flagged sets are disjoint, and
\begin{equation}
 |E_C|\le R_{E,H}n_C,\qquad
 E_C\text{ has a connected hull of at most }R_{E,H}n_C\text{ sites},
 \qquad R_{E,H}\le Ct_H^{-2}.
 \label{cmp:edit-size}
\end{equation}
One can replace the spins in any selected \(E_C\), leaving all other
spins and the region's outer boundary unchanged, so that every true
rough bond of \(D_C\) disappears and no true rough bond is created.
All choices can be made measurably.
\end{lemma}

\begin{proof}
Partition each coordinate circle into integer intervals of lengths
between \(\ell=\lceil C_2/t_H\rceil\) and \(2\ell\).  This is possible
once \(M\) exceeds a fixed multiple of \(\ell\).  Use the resulting
rectangular mesh.  Mark every box meeting a bond of \(\widehat D\).
Take the union of their closed two-box neighborhoods, and form
components by joining boxes that meet even at a vertex.  Distinct
components are disjoint closed mesh regions.  Keep just the components
containing a bond of \(D\).

Let \(E_C\) be the lattice sites in one retained region, and let \(n_C\)
count all originally flagged bonds whose marked boxes generate that
region.  A flagged bond marks only a bounded number of boxes, and each
marked box adds only a bounded number of boxes to its neighborhood.
Hence the number of boxes is at most \(Cn_C\), proving
\eqref{cmp:edit-size}.  The regions have a connected lattice hull with
the same bound.  Every flagged bond is separated from the boundary of
its component by a mesh collar.  Thus boundary bonds are unflagged.

Retain the original spins on all boundary vertices and edges of the
region.  Interpolate each such edge by short geodesic arcs between
successive lattice spins.  Its speed in physical lattice coordinates
is at most \(Ct_H/C_1\).  At every other mesh vertex choose the fixed
spin \(\mathbf1\).  On the other mesh sides join the chosen endpoint
spins by a fixed measurably selected shortest arc.  Use exactly the same
arc for the two cells sharing a side.

After rescaling a cell to the unit square, its boundary map has
Lipschitz constant at most \(C(1+C_2/C_1)\).  Apply
Lemma~\ref{cmp:filling} in every cell.  Choose \(C_2\) larger than the
resulting fixed filling constant for \(C_2/C_1\le1\), and then choose
\(C_1\) large enough.  The Lipschitz constant in lattice units is less
than \(t_H\).  Sampling the fillings at lattice sites therefore gives
new bond lengths below \(t_H\).

The fillings agree on shared sides.  Boundary spins were preserved,
so no bond crossing out of a region becomes rough.  The construction
also treats regions with holes.  If a region wraps around the torus,
the same shared-side rule is used periodically; if it fills the torus,
one may simply replace all its spins by \(\mathbf1\).
\end{proof}

Fix the regions and the modified values using the original
configuration, once and for all.  Write \({\cal I}\) for the retained
regions and \(q(I)\), \(I\subset{\cal I}\), for the configuration in
which precisely the regions outside \(I\) have been modified.  Then
\begin{equation}
 D(q(I))=\bigsqcup_{C\in I}D_C.
 \label{cmp:inventory-after-edit}
\end{equation}
The modification need not remove the wider flags at threshold
\(t_H/C_1\).  We never recluster a modified configuration; its geometry,
its \(n_C\)'s, and its inventories are those fixed from the original one.

Put \(\Xi_i(I)=\Xi_i(q(I))\).  Since \(\rho_i>0\), these are complete
positive sums.  Equations \eqref{cmp:density-finite-energy} and
\eqref{cmp:exponent-finite-energy} give the crucial denominator estimate:
\begin{equation}
 0<\frac{\Xi_i(I\setminus B)}{\Xi_i(I)}
 \le
 \exp\!\left(C_{3,H}\sum_{C\in B}n_C\right),
 \qquad
 C_{3,H}=(B_{\rho,H}+B_{X,H})R_{E,H}.
 \label{cmp:complete-ratio}
\end{equation}
Indeed \(\log\Xi_i=\log\rho_i-X_i\), and at most
\(R_{E,H}\sum_{C\in B}n_C\) sites have changed.

\subsection{Exact identities for the signed covering sums}

At a fixed active set \(I\), take a full original covering family and
join two of its labels whenever their supports meet the same active
region \(E_C\), \(C\in I\).  Call a connected group a macrolabel \(U\),
and give it support
\begin{equation}
 S_U=\bigcup_{\lambda\in U}P_\lambda
       \ \cup\
       \bigcup_{\substack{C\in I:\\
                    E_C\cap P_\lambda\ne\varnothing
                    \text{ for some }\lambda\in U}}E_C.
 \label{cmp:macro-support}
\end{equation}
Its weight is
\[
 w_i(U;q(I))=\prod_{\lambda\in U}k_{i,\lambda}(q(I)).
 \label{cmp:macro-weight}
\]
Every active region it meets is covered in full by this macrolabel:
all labels supplying its inventory meet that region and so belong to
the same group.  Different macrolabels have disjoint enlarged supports.
Conversely, compatible macrolabels covering all active inventories
recover exactly the original full covering families.

For a site set \(S\) disjoint from the active editable regions, let
\(\Xi_i(I;S)\) be the full-inventory sum restricted to macrolabels whose
supports avoid \(S\).  This restricted sum may have either sign.  If
\({\cal A}\) is a compatible selected macrolabel family, denote its
combined support by \(S_{\cal A}\), and the active regions it covers by
\(B_{\cal A}\).

\begin{lemma}[Removal and inclusion--exclusion]
\label{cmp:removal-identities}
The background left after selecting \({\cal A}\) has total weight
\begin{equation}
 \Xi_i(I\setminus B_{\cal A};S_{\cal A}).
 \label{cmp:background-identity}
\end{equation}
Consequently
\begin{align}
 \Xi_i(I;S)
 &=
 \sum_{\substack{{\cal A}\ {\rm compatible}\\
                   S_U\cap S\ne\varnothing\ \forall U\in{\cal A}}}
 (-1)^{|{\cal A}|}
 \left[\prod_{U\in{\cal A}}w_i(U;q(I))\right]
 \Xi_i(I\setminus B_{\cal A};S_{\cal A}),                 \label{cmp:IE}\\
 \Xi_2(I)-\Xi_1(I)
 &=
 \sum_{\substack{{\cal A}\ {\rm compatible}\\{\cal A}\ne\varnothing}}
 \left[\prod_{U\in{\cal A}}
       \bigl(w_2(U;q(I))-w_1(U;q(I))\bigr)\right]
 \Xi_1(I\setminus B_{\cal A};S_{\cal A}).                 \label{cmp:difference-identity}
\end{align}
The empty term in \eqref{cmp:IE} is \(\Xi_i(I)\).
\end{lemma}

\begin{proof}
The configurations \(q(I)\) and \(q(I\setminus B_{\cal A})\) differ
only inside \(S_{\cal A}\).  Every remaining weight and every
eligibility indicator reads a support avoiding this set, so its value
is unchanged.  A remaining active region is disjoint from
\(S_{\cal A}\), by the definition of \(B_{\cal A}\).  Removing the
selected regions therefore does not change the grouping of background
labels: background supports avoided these regions already.  Its
inventory requirement loses exactly the inventories in
\(B_{\cal A}\).  This proves the weight-preserving bijection
\eqref{cmp:background-identity}.

In every full covering family, expand the indicator that all its
macrolabels avoid \(S\) as
\[
 \prod_U\bigl(1-\mathbf1_{\{S_U\cap S\ne\varnothing\}}\bigr).
\]
After selecting the factors carrying the minus sign, apply
\eqref{cmp:background-identity}; this gives \eqref{cmp:IE}.  For
\eqref{cmp:difference-identity}, expand each product using
\(w_2=w_1+(w_2-w_1)\) and subtract the all-\(w_1\) term.

One may first perform these algebraic manipulations with finite label
lists.  A full family has at most \(|D(q(I))|\) original labels, and
\eqref{cmp:absolute-cover} provides absolute convergence of all the
unnormalized sums.  The identities therefore pass to the countable
lists.  Division by complete sums is performed only after this limit;
no positivity of a truncated or restricted sum is required.
\end{proof}

Introduce the nonnegative envelopes
\[
 p_\lambda=\|k_{1,\lambda}\|_\infty+
                         \|k_{2,\lambda}\|_\infty,\qquad
 d_\lambda=\|k_{2,\lambda}-k_{1,\lambda}\|_\infty
\]
and the restricted ratio
\[
 R_i(I,S)=\frac{|\Xi_i(I;S)|}{\Xi_i(I)}.
\]
Equations \eqref{cmp:complete-ratio} and \eqref{cmp:IE} give
\begin{equation}
 R_i(I,S)\le1+
 \sum_{\substack{{\cal A}\ne\varnothing,\ {\rm compatible}\\
                         S_U\cap S\ne\varnothing\ \forall U\in{\cal A}}}
 \left[\prod_{\lambda\in\bigcup{\cal A}}p_\lambda\right]
 e^{C_{3,H}\sum_{C\in B_{\cal A}}n_C}
 R_i(I\setminus B_{\cal A},S_{\cal A}).
 \label{cmp:ratio-recursion}
\end{equation}
A telescoping of each finite product also gives
\begin{equation}
 |w_2(U)-w_1(U)|
 \le\sum_{\lambda\in U}d_\lambda
                   \prod_{\substack{\kappa\in U\\\kappa\ne\lambda}}
                          p_\kappa .
 \label{cmp:marked-macro}
\end{equation}
Thus \eqref{cmp:difference-identity} starts the same recursion with a
nonempty first generation and one marked, differenced original label
in each of its macrolabels.  Every generation consumes at least one
active region.  The recursion is consequently finite.

\subsection{The common exceptional set}

Let \(C_{4,H}>0\) be a polynomial in \(H\), to be fixed in the next
subsection.  Define \(G\) by the following requirement.  Take any finite
collection of distinct original label nodes and distinct editable-region
nodes whose physical supports form a connected intersection graph.
Write
\[
 s=\sum_{\lambda}s_\lambda,\qquad n=\sum_C n_C.
\]
Whenever \(n>0\), require
\begin{equation}
 (C_{3,H}+C_{4,H})n\le \frac A4s+\zeta_HK.
 \label{cmp:good-set}
\end{equation}
This quantification uses the whole common label set, not just labels
eligible at the original configuration.  The recursion can evaluate a
label at a different \(q(I)\), and the same geometric inequality must
remain available there.  All these collections are countable, so \(G\)
is measurable.

\begin{lemma}[Probability of the exceptional set]
\label{cmp:exceptional}
For every polynomial choice of \(C_{4,H}\), sufficiently large \(H\)
admits an \(a_H>0\) such that
\begin{equation}
 \mu_i(G^c)\le2v e^{-a_HK},\qquad i=1,2 .
 \label{cmp:exceptional-bound}
\end{equation}
\end{lemma}

\begin{proof}
A violation of \eqref{cmp:good-set} implies
\[
 n>\frac{\zeta_HK}{C_{3,H}+C_{4,H}},\qquad
 s<\frac{4(C_{3,H}+C_{4,H})}{A}n.
\]
The connected hull bounds on labels and editable regions give a
connected graph containing these \(n\) distinct originally flagged
bonds, with at most \(B_Hn\) vertices and edges, for a polynomial
\(B_H\).  A traversal of this graph visits all the bonds with total
length at most \(CB_Hn\).

We count the ordered flagged bonds in such a traversal, rather than
every intervening lattice step.  There are at most \(2v\) choices for
the first bond.  At distance \(\ell\), there are at most
\(C(\ell+1)\) possible next bonds, including their orientations.
If \(Q\le CB_Hn\), the generating identity
\[
 \sum_{\ell\ge0}(\ell+1)z^\ell=(1-z)^{-2}
\]
therefore gives
\[
 \sum_{\ell_1+\cdots+\ell_n\le Q}
              \prod_{j=1}^n(\ell_j+1)
 =\binom{Q+2n}{2n}
 \le [C(1+B_H)^2]^n .
\]
This also bounds the number of ordered distinct-bond lists; allowing
repetitions in the count only increases it.  On a torus one may lift
the traversal before projecting its recorded bonds, which gives the
same upper bound.  Thus the number of possible flagged sets forced by
a violation is at most
\begin{equation}
 v[C(1+B_H)^2]^n .
 \label{cmp:flag-entropy}
\end{equation}
This estimate costs \(O(\log H)\) per flag, not \(O(B_H)\) per flag.

Set \(Q_H=C(1+B_H)^2e^{-\sigma_H}\).  Since \(B_H\) is polynomial and
\(\sigma_H=cHt_H^2\), assumption (C4) gives \(Q_H<1/2\) for large \(H\).
Lemma~\ref{cmp:rarity} and \eqref{cmp:flag-entropy} imply, under either law,
\[
 \mu_i(G^c)
 \le v\sum_{n>\zeta_HK/(C_{3,H}+C_{4,H})}Q_H^n
 \le2v e^{-a_HK},\qquad
 a_H=\frac{\zeta_H\log(1/Q_H)}{C_{3,H}+C_{4,H}}>0 .
\]
The possibly infinite number of labels adds no factor: every violating
collection already forces one of the flagged sets counted above.
\end{proof}

\subsection{A tree bound for the signed recursion}

We now prove the normalized estimate, including its combinatorial
multiplicity.

\begin{lemma}[Normalized covering comparison]
\label{cmp:forest}
For a sufficiently large polynomial \(C_{4,H}\) and sufficiently large
\(H\), the set \(G\) above satisfies
\begin{equation}
 \frac{|\Xi_2(q)-\Xi_1(q)|}{\Xi_1(q)}
 \le \exp(D_Hv\delta_K e^{\zeta_HK})-1 .
 \label{cmp:forest-conclusion}
\end{equation}
The same estimate holds with denominator \(\Xi_2(q)\).
\end{lemma}

\begin{proof}
Iterate \eqref{cmp:ratio-recursion} in
\eqref{cmp:difference-identity}, using \eqref{cmp:marked-macro}.  A term
is specified by successive nonempty compatible macrolabel families;
the recursion ends when the term \(1\) is chosen.  Every macrolabel
after the first generation meets the combined support of the preceding
generation.

Here is a faithful encoding of these terms by forests.  Inside each
macrolabel choose a spanning tree of its original label nodes and
editable-region nodes, using a fixed deterministic ordering to make the
choice.  For each later macrolabel choose, again deterministically, one
support intersection joining it to a macrolabel of the preceding
generation.  Color an edge according to whether it is internal to a
macrolabel or joins two generations.  Root each resulting tree at the
marked original label in its initial macrolabel.

No node is repeated.  An editable region is consumed the first time it
appears.  An original label has a fixed nonempty inventory; after it
appears, the regions containing that inventory have been removed, so
it cannot appear again in an eligible subsequent family.  Within one
generation supports are disjoint.

The encoding loses no multiplicity. Contract the internal edges to
recover each macrolabel, including its original labels and editable
regions. The depth of the resulting vertex in its rooted tree
recovers its generation; vertices at the same depth, across all trees,
recover the compatible family selected at that stage. The marked roots
recover the choices in the product telescoping. Finally, remove the
regions in all earlier generations to recover the active set \(I\) at
that stage and hence the configuration \(q(I)\) at which every weight
was evaluated. These data specify the original recursive term.
Thus two different terms cannot produce the same forest with its node
labels, marks and two edge colors. No ordering of siblings, or other
choice of a recursive history, remains to be counted.

Each tree is a connected collection of distinct nodes of the kind used
in \eqref{cmp:good-set}.  Its accumulated denominator factor is therefore
bounded by
\begin{equation}
 e^{C_{3,H}n}
 \le e^{\zeta_HK}
       \prod_{\lambda\ {\rm in\ the\ tree}}e^{As_\lambda/4}
       \prod_{C\ {\rm in\ the\ tree}}e^{-C_{4,H}n_C}.
 \label{cmp:tree-charge}
\end{equation}
The factor \(e^{\zeta_HK}\) occurs once per rooted tree.

To sum the remaining forests, give a gas node size \(s_\lambda\) and
an edit node size \(|E_C|\).  Their unmarked weights are
\[
 p_\lambda e^{As_\lambda/4},\qquad e^{-C_{4,H}n_C},
\]
respectively.  A marked gas root instead has weight
\(d_\lambda e^{As_\lambda/4}\).  Put \(R=\max(1,R_H)\) and choose
\begin{equation}
 C_{4,H}\ge2R_{E,H}+\log(16R)+1 .
 \label{cmp:C4}
\end{equation}
Since the editable sets are disjoint and \(n_C\ge1\),
\begin{equation}
 \sup_x\sum_{C:x\in E_C}
             e^{-C_{4,H}n_C}e^{2|E_C|}
       \le\frac1{16R}.
 \label{cmp:edit-hit}
\end{equation}
By \eqref{cmp:activity}, \(A/4+2\le A\), and the rapid decay of \(w_H\),
\begin{equation}
 \sup_x\sum_{\lambda:x\in P_\lambda}
           p_\lambda e^{As_\lambda/4}e^{2s_\lambda}
       \le2w_H\le\frac1{16R}.
 \label{cmp:gas-hit}
\end{equation}

A node of size \(a\) has support at most \(Ra\).  Summing a potential
child over one of its intersection sites and over the two edge colors,
equations \eqref{cmp:edit-hit}--\eqref{cmp:gas-hit} bound the child sum,
even with an exponential allowance for its descendants, by \(a/4\).
More explicitly, induction on maximum tree depth bounds the sum of
descendants below a root of size \(a\) by \(e^a\).  For the induction
step, each child of size \(b\) contributes at most its weight times
\(e^b\).  The sum over unordered distinct children is at most the
exponential of their one-child sum: sum ordered \(r\)-tuples with
factor \(1/r!\) and then drop distinctness.  The result is at most
\(e^{a/4}\le e^a\).  Positivity of these majorants permits the limit in
the maximum depth.

The marked-root sum is bounded by
\begin{equation}
 \sum_\lambda d_\lambda e^{As_\lambda/4}e^{s_\lambda}
       \le D_Hv\delta_K.
 \label{cmp:root-hit}
\end{equation}
Here summing the per-site bound can only overcount labels.  Actual
forest components have distinct marked roots.  Drop all disjointness
conditions between their descendant trees and sum unordered root sets.
Equations \eqref{cmp:tree-charge} and \eqref{cmp:root-hit} give
\[
 \sum_{\text{nonempty forests}}\text{majorant}
 \le \sum_{r\ge1}
       \frac{(D_Hv\delta_K e^{\zeta_HK})^r}{r!},
\]
which is \eqref{cmp:forest-conclusion}.

Every denominator used above was a complete positive \(\Xi_i(I)\);
every restricted signed sum entered only through its absolute value.
The argument is symmetric in \(1,2\), proving the last assertion.
\end{proof}

\subsection{From relative densities to partition functions}

\begin{lemma}[Normalization on a common set]
\label{cmp:normalization}
Let \(f_1,f_2>0\) be integrable, with \(0<Z_i=\int f_i\,\dd\nu<\infty\)
and \(\pi_i=f_i\nu/Z_i\).  If
\[
 \pi_i(G^c)\le\varepsilon_i<1,\qquad
 \left|\log\frac{f_2}{f_1}\right|\le\eta\quad\hbox{on }G,
\]
then
\begin{equation}
 e^{-\eta}(1-\varepsilon_1)
       \le\frac{Z_2}{Z_1}
       \le\frac{e^\eta}{1-\varepsilon_2}.
 \label{cmp:normalization-bound}
\end{equation}
If \(\varepsilon_i\le\varepsilon\le1/2\), then
\(\left|\log(Z_2/Z_1)\right|\le\eta+2\varepsilon\).
\end{lemma}

\begin{proof}
The relative density bound compares the two integrals over \(G\)
by factors \(e^{\pm\eta}\).  Use the exact identity
\[
 \frac{Z_2}{Z_1}
   =\frac{\int_G f_2\,\dd\nu}{\int_G f_1\,\dd\nu}
                       \frac{\pi_1(G)}{\pi_2(G)}.
\]
The last assertion follows from
\(-\log(1-\varepsilon)\le2\varepsilon\).
\end{proof}

\begin{proof}[Proof of Theorem~\ref{thm:relative-comparison}]
Choose \(C_{4,H}\) as in \eqref{cmp:C4}.  Lemma~\ref{cmp:exceptional}
gives the common exceptional-set estimate, and
Lemma~\ref{cmp:forest} gives \eqref{cmp:relative-cover}.

Applying the latter estimate in both directions yields
\[
 e^{-u}\le\frac{\Xi_2(q)}{\Xi_1(q)}\le e^u,\qquad
 u=D_Hv\delta_K e^{\zeta_HK}.
\]
This step does not require \(u\) to be small.
Adding \eqref{cmp:exponent-difference} proves
\eqref{cmp:relative-density}.

Let
\[
 b_H=\min(a_H,\gamma_H-\zeta_H)>0
\]
and choose \(0<c_H<b_H/2\).  If \(v\le e^{c_HK}\), both the exceptional
probability and the relative density error are at most
\(C_He^{-(b_H-c_H)K}\).  For \(K\) sufficiently large the former is
at most \(1/2\), so Lemma~\ref{cmp:normalization} proves
\eqref{cmp:partition-conclusion}, with any fixed
\(0<d_H\le b_H-c_H\) after adjusting \(C_H\).
\end{proof}

\subsection{Verification for the renormalized \(O(4)\) densities}

We now check the assumptions used in the theorem, including those that
come from positivity of the original model rather than from bounds on
its signed expansion.

At the retained scale, Theorem~\ref{thm:rg-closure} gives
\begin{equation}
 \rho_i=e^{X_i}\Xi_i,\qquad
 X_i=-h{\cal E}_i-{\cal P}_i-{\cal L}_i,\qquad
 H/2\le h\le2H.
 \label{cmp:RG-density}
\end{equation}
The two densities are written using the prescribed bulk scalar
normalization.  The kinetic form is
\begin{equation}
 {\cal E}_i(q)=\sum_e S_t(r_e)
           +\frac12\sum_e m_e(q)|(\ell_i\nabla q)_e|^2,
 \qquad
 S_t(r)=
 \begin{cases}
       c_0r^2/2,&r\le t,\\
       d_0tr,&r>t,
 \end{cases}
 \label{cmp:kinetic-form}
\end{equation}
where \(0<d_0<c_0/16\).  In this formula \(\nabla q\) denotes the field of
nearest-neighbor spin differences.  The free kernels \(\ell_i\) have
the stated uniform exponentially weighted row and column bounds.
Every mask and every other condition used by a term is part of its
complete support.

\paragraph{Support and difference norms.}
The complete-cover convention gives a connected hull with at most
\(C_L{\cal M}_H^2s_\lambda\) sites, where
\({\cal M}_H=\lceil(\log H)^2\rceil\).  Thus \(R_H\) is polynomial.
The covering bound is
\[
 w_H=\exp[-p_H^{1/4}],
\]
which decreases faster than every inverse power because \(P_0>4\).
The compatible lists and the equal-endpoint estimate of
Theorem~\ref{rg:matching} give
\eqref{cmp:activity-difference} and \eqref{cmp:exponent-difference}
with \(\delta_K\le C_He^{-\gamma_HK}\).  Any fixed polynomial in \(K\)
in that estimate is absorbed by decreasing \(\gamma_H\).
The constants are uniform in the longer cutoff depth and in the
admitted terminal coupling.

\paragraph{Strict positivity and finite energy of the full density.}
The last block observation has conditional density
\[
 Q_b(V\mid u)=z_{ab}^{-1}e^{abV\cdot u},\qquad |u|=1.
\]
In the small-mean branch it is a positive mixture of these densities
over constituent-spin directions.  In either case,
\begin{equation}
 e^{-2ab}\le\frac{Q_b(V\mid\cdot)}{Q_b(\widetilde V\mid\cdot)}
                     \le e^{2ab}.
 \label{cmp:kernel-ratio}
\end{equation}
For a mixture the bound follows term by term before summation.
The observation variables in different cells are conditionally
independent.  Changing \(n\) output spins therefore changes the product
kernel by a factor between \(e^{-2abn}\) and \(e^{2abn}\); integration
against any positive input density preserves this bound.

The incoming coupling in the last step is at most
\(2(H+\gamma_0)\), where \(\gamma_0=(\log L)/\pi\).  Hence
\[
 B_{\rho,H}=4a(H+\gamma_0)
\]
works in \eqref{cmp:density-finite-energy}, uniformly in all preceding
steps.  It also follows directly that the output density is strictly
positive.  Scalar extraction does not affect the ratio bound.

\paragraph{Reflection positivity.}
At a retained seam, the fine reflection exchanges the two disjoint
halves of the block partition.  The symmetric block weights, the
normalized-mean rule, and the constituent-spin fallback commute with
that reflection.  For a function \(F\) of output spins on one side,
let \({\cal Q}F\) be its conditional expectation given the input spins
on that side.  Conditional independence and reflection covariance give
\[
 \mu_{\rm out}((\theta\overline F)F)
   =\mu_{\rm in}
       ((\theta\overline{{\cal Q}F})\,{\cal Q}F)\ge0.
\]
Thus reflection positivity passes through each blocking step.
Translation and reflection compatibility on the retained lattice pass
through the same kernels.  Notice that conditional independence is
used together with reflection covariance.

\paragraph{Stability.}
The second term of \eqref{cmp:kinetic-form} is nonnegative and
\[
 h{\cal E}_i(q)\ge
           \frac{d_0H}{2}\sum_e\min(r_e^2,t_H^2).
\]
The seven canonical slots have degrees \(4,2,5,3,6,4,2\) and prefactors
\(h,1,h,1,h,1,h^{-1}\).  On a slot's regular mask, a displacement
along a recorded path of length \(u\) is at most \(ut_H\).
The exponential path norm absorbs every fixed power of \(u\), so
\begin{align*}
 \|{\cal P}_i\|_\infty
 &\le C_Lv\bigl(Ht_H^4+t_H^2+Ht_H^5+t_H^3
                         +Ht_H^6+t_H^4+H^{-1}t_H^2\bigr)\\
 &\le C_Lv(Ht_H^4+t_H^2).
\end{align*}
The regular-error norm gives
\(\|{\cal L}_i\|_\infty\le vH^{-2.05}\).  Thus
\eqref{cmp:stability} holds with
\[
 c_*=\frac{d_0}{2},\qquad
 \epsilon_H=C_L(Ht_H^4+t_H^2+H^{-2.05})=o(Ht_H^2).
\]
On the cap \({\cal C}_H\), all bonds are \(O(H^{-1/2})\), all masks
are on, and the row bound for \(\ell_i\) gives
\(h{\cal E}_i\le C_Lv\).  The same canonical and error estimates yield
\(X_i\ge-C_Lv\) there.  Since \(Ht_H^2\to\infty\), this is
\eqref{cmp:cap}.

\paragraph{Finite energy of the regular exponent.}
Changing \(n\) spins changes only \(O(n)\) direct bonds.  The row and
column bounds for \(\ell_i\) imply
\[
 \sum_e|(\ell_i\nabla(q-\widetilde q))_e|\le C_L n.
\]
A kinetic mask can change only if its read neighborhood meets a changed
bond.  Its radius is \(O(\log(2/t_H))\), so there are at most
\(C n\log^2(2/t_H)\) such masks.  Uniform boundedness of the row values
therefore bounds the kinetic variation by
\[
 C_LH[1+\log^2(2/t_H)]\,n.
\]
For the other terms, complete supports imply that a change is possible
only in terms whose supports hit a changed site.  Their exponential
path and cover norms bound this hit sum by a fixed polynomial in
\({\cal M}_H\) times their anchored norms: sum over possible anchors
in a connected hull, and absorb its fixed size powers with the spare
exponential weight.  It follows that
\[
 B_{X,H}\le C_LH[1+{\cal M}_H+\log(2/t_H)]^C,
\]
which proves \eqref{cmp:exponent-finite-energy}.

Finally \(t_H\to0\), \(t_H^{-1}\) is polynomially bounded, and
\(Ht_H^2=p_H^2\) dominates \((\log H)^2\).  This completes the
verification of (C1)--(C4).  Increasing the fixed minimum period to
accommodate the boxes in Lemma~\ref{cmp:editing} is harmless: it depends
only on \(H\), not on either cutoff depth.

